%% chapters/20-bbp-transition.tex
\chapter{The BBP Phase Transition}
\label{ch:bbp}

\whereWeAre{%
The BBP (Baik--Ben Arous--P\'ech\'e) phase transition is the precise mathematical statement of when a spike escapes the Marchenko--Pastur bulk. There is a critical signal strength below which the spike is invisible (eigenvalue stays at the bulk edge) and above which it becomes detectable (eigenvalue moves outside the bulk). This is a phase transition in the literal mathematical sense.%
}

\PLACEHOLDER{Mostly new chapter. Reference: Yale Stat 615 Lecture 9, original BBP paper.}

\section{Sub-critical and super-critical regimes}
\PLACEHOLDER{Critical strength $\theta^* = \sqrt c$. Below: spike eigenvalue stays at $\lambda_+$. Above: spike eigenvalue at $\theta + c\theta/(\theta - 1)$ or similar formula.}

\section{Three pictures of the transition}
\PLACEHOLDER{Sub-critical histogram, super-critical histogram, $\theta$ sweep showing the phase transition.}

\section{Detectability and PCA}
\PLACEHOLDER{Why this matters: PCA is only useful above the BBP threshold. Below threshold, the top sample eigenvector is essentially random.}

\section{Simulation: BBP transition sweep}
\PLACEHOLDER{Colab box: parameter sweep over $\theta$ with histogram of top sample eigenvalue.}

\section{Brief warm-ups}
\PLACEHOLDER{3--5 short questions.}

\chapterRecap{%
\begin{itemize}
  \item \PLACEHOLDER{BBP gives a sharp threshold for spike detection.}
  \item \PLACEHOLDER{Below threshold: spike invisible.}
  \item \PLACEHOLDER{Above threshold: spike detectable, eigenvalue moves outside the bulk.}
\end{itemize}%
}

\section*{Exercises}
\PLACEHOLDER{8--15 longer exercises.}
